\section{Introduction}

Static bug analyzers play a crucial role in ensuring software security, reliability, and maintainability. They help identify potential bugs early in the development process by scanning the codebase without extensive program execution~\cite{cai2023place,cai2022peahen,cai2021canary,shi2020conquering,shi2018pinpoint,sui2012static,vassallo2020developers,yan2018spatio}. Popular analyzers like GitHub CodeQL~\cite{codeql}, Facebook Infer~\cite{infer}, and CppCheck~\cite{cppcheck} are widely used for industrial bug detection. Their analysis of many bug types (\eg{} null pointer dereference and use-after-free) typically follows a standardized approach: identifying sources (\ie{} potential origins of issues) and sinks (\ie{} locations where issues manifest) based on predefined rules, then checking whether a feasible execution path exists from sources to sinks. If such a path is found, the tools flag a potential bug.

However, in complex software systems, sources and sinks often reside in different functions. In cases where function call chains between a source and a sink are lengthy and involve multiple conditional branches and complex data dependencies, traditional static analyzers struggle to comprehensively validate path feasibility between sources and sinks, resulting in high false positive rates. For instance, our manual annotation of sampled bug reports reveals that popular tools produce an extremely high rate (e.g., over 90\%) of false positives for use-after-free bugs across large-scale software such as Linux Kernel. This significantly burdens developers due to the large efforts in manually reviewing the false alarms reported by these tools. 

\textbf{Existing Techniques.} As Large Language Models (LLMs) are increasingly applied to code-related tasks~\cite{du2024evaluating,liu2024large,nam2024using, du2024vulragenhancingllmbasedvulnerability}, researchers have leveraged their advanced program understanding capabilities to reduce false positives generated by static analyzers. However, existing LLM-based approaches have limited effectiveness in handling path feasibility analysis within large, dependency-rich codebases. Specifically, LLM4SA~\cite{llm4sa} feeds entire source-to-sink call chains along with bug reports into LLMs at once to identify false positives, without conducting fine-grained constraint analysis and relying entirely on the LLM's ability to assess complex path feasibility. LLMDFA~\cite{llmdfa} performs data flow analysis for each function along all potential execution paths between sources and sinks, which results in significant time and computational cost when applied to large-scale projects with complex call dependencies.
While symbolic execution can perform path feasibility analysis, it suffers from path explosion when traversing all possible paths, particularly when dealing with multiple conditional branches, loops, nested function invocations, or compound data types~\cite{10.1145/3182657}. Therefore, it cannot scale to large and complex software systems in practice. 

\textbf{Our approach. }
To address these limitations, we propose \app{}, an LLM-powered path feasibility analysis framework designed to minimize false positives in static bug detection for complex software projects. 
Specifically, Given a function call trace and a target variable with its source and sink, \app{} iteratively performs two-phase analysis for each function: (1) \textit{Feasibility Constraints Extraction} first identifies critical conditional expressions impacting sink reachability; then employs LLM agents to conduct symbolic range reasoning for variables and function calls within each expression. (2) \textit{Constraints Solving} employs constraint solver to iteratively solve the SMT query scripts generated by LLM. 
Overall, \app{} offers the following advantages compared to traditional symbolic execution and existing LLM-based methods: 

\begin{itemize}[itemsep=2pt,topsep=0pt,parsep=0pt]
\item \textit{Precise complex constraint cascade analysis.} Breaking down inter-procedural feasibility analysis into localized symbolic range reasoning, \app{} amplifies LLMs' program comprehension strength on focused code segments~\cite{li2023loogle}, thereby significantly enhancing reasoning effectiveness.

\item \textit{Efficient complex context analysis.} Integrating prioritized \constraint{} extraction method with agent-guided planning, \app{} effectively mitigates path explosion in  sequentially traversing deeply-nested contextual functions.
\end{itemize}


\textbf{Evaluation.}
We manually construct a new dataset, SAFP-BENCH-C, for a more realistic evaluation, which takes approximately 100 person-hours.  It contains 364 warnings generated by three advanced static analyzers, covering three important bug types (i.e., null pointer dereference, buffer overflow, and use after free) in large-scale software projects, including the Linux Kernel (with over 18 million lines of code), OpenSSL (with over 300K lines of code), and Libav (with over 600K lines of code). Our evaluation shows that \app{} precisely filters out 72\% to 96\%) false positives reported during static bug detection, significantly outperforming all the baselines by 41.1\% - 105.7\% improvements; meanwhile \app{} only misses 3 real bugs of 45 true positives. 
Furthermore, \app{} shows consistent effectiveness across four state-of-the-art LLMs, indicating the generalizability of our approach. Ablation studies further reveal that \app{} outperforms two ablation variants, underscoring the key contribution of LLM agent based self-planning context analysis and iterative constraint solving proposed in our framework.

This paper makes the following contributions:

\begin{itemize}[itemsep=2pt,topsep=0pt,parsep=0pt]

\item \textbf{A novel technique} that effectively enhances  complex inter-procedural feasibility analysis to minimize false positives in bug detection.

\item \textbf{A new dataset} designed for a more realistic evaluation of false positive reduction, with approximately 100 person-hours dedicated to manual construction.

\item \textbf{An extensive evaluation} demonstrating the effectiveness of \app{} in filtering out false positives and identifying real bugs. 
\end{itemize}

